\subsection*{Abstract}
\noindent In a Korean \captive{} a listed parent moves a business into a wholly owned subsidiary; whether it is later kept, sold or listed is decided by the controlling owner without a shareholder vote, and any loss to the minority comes at the disposal. The literature on controlling-owner expropriation infers intention from the expropriating outcome; we measure beforehand what the filing disclosed about the risk that the subsidiary would not be kept. Two language models read the division filings of \NScored{} \captive{}s of 2001--2026 and forecast whether the parent will still own the subsidiary three to five years later. For the \NGateOne{} labelled divisions the forecast predicts the realised disposal (AUC \AUCPooledLLM{}) beyond the filing's financials and the purpose categories and mention flags a hand coder would record, survives the removal of every name and date, and doubles the hazard of disposal per standard deviation. The announcement return rises with the forecast, the following year's return falls with it against market and portfolio benchmarks, and the loss is concentrated where the subsidiary is disposed of within a year. Smaller controlling blocks go with stronger forecasts and more disposals. The 2022 reform's measures left separable traces: against equity spin-offs it raised the announcement reaction by \DiDCxP{} percentage points while the mean forecast moved from \ScorePre{} to \ScorePost{} ($n$ = \ProtPostN{}); the appraisal right is priced as a put, imprecisely; and the forecast is lower only where the filing binds the parent to keep the subsidiary.
\par\medskip{\singlespacing\noindent\textbf{Keywords:} captive spin-offs; controlling shareholders; minority shareholder protection; corporate disclosure regulation; appraisal rights; business groups; Korea
\par\noindent\textbf{JEL classification:} G32; G34; G38; M41; M48\par}
